% IEEE-format plain-language ("simplified") version of the MIC paper.
% Prose is the accessible/simplified narrative; tables (no figures) are
% reproduced from mic-paper-v9-ieee-tmi.tex; references are the cited
% subset of the v9 bibliography.
% Compile: pdflatex -> pdflatex
\documentclass[journal]{IEEEtran}

% ---- Math / fonts ----
\usepackage{amsmath}
\usepackage{amssymb}
\usepackage{newtxtext,newtxmath}

% ---- Graphics / tables ----
\usepackage{graphicx}
\usepackage{booktabs}
\usepackage{array}
\usepackage{multirow}

% ---- Utilities ----
\usepackage{xcolor}
\usepackage{url}
\usepackage{hyperref}
\usepackage{microtype}
\usepackage{cite}

\hypersetup{
  colorlinks=true,
  linkcolor=blue!70!black,
  citecolor=green!50!black,
  urlcolor=blue!70!black,
}

% --------------------------------------------------
\title{A 16-Bit-Native Compression Pipeline for\\
       Lossless Medical Images\\
       \large(Plain-Language Edition)}

\author{Kuldeep~Singh%
\thanks{K.\ Singh is with Innovaccer Inc.
E-mail: kuldeep.singh@innovaccer.com.
ORCID: 0009-0004-8476-0118.
Code: \url{https://github.com/pappuks/medical-image-codec}.}%
\thanks{This is a plain-language edition of the MIC paper. It keeps the main
ideas, key numbers, and reasoning, but trades the formal equations and
implementation minutiae for explanations a software engineer can follow.
Data tables are reproduced from the full manuscript; figures are omitted.}%
}

% --------------------------------------------------
\begin{document}
\maketitle

% ============================================================
\begin{abstract}
Medical images must be stored and transmitted losslessly, since a radiologist
must see exactly what the scanner produced, and they are ``deep,'' carrying
10--16 bits per pixel rather than the 8 bits of an ordinary photo. That deep
format is where most compression tools perform poorly, because nearly all of
them were built for 8-bit data. This paper describes MIC (Medical Image Codec),
a lossless codec that works directly on 16-bit data instead of splitting each
pixel into two bytes. MIC chains three simple stages, a spatial predictor, a
16-bit run-length encoder, and a large-alphabet entropy coder called Finite
State Entropy (FSE, a fast table-driven form of asymmetric numeral systems),
and adds a multi-state decoder that keeps the CPU more fully utilized to speed
up decompression. On 39 real DICOM images across eleven modalities, MIC
compresses about 10\% better on average (better on 34 of the 39 images) than a
strong general-purpose baseline (Delta\,+\,Zstandard), reaches a geometric-mean
compression ratio of 3.36$\times$ (within about 2\% of High-Throughput
JPEG~2000 at 3.42$\times$ and about 88\% of JPEG-LS at 3.84$\times$), and is the
fastest decoder of the four tested codecs on all 39 images on ARM64 and 36 of 39
on AMD64. A 20\,KB pure-JavaScript decoder lets the format be decoded directly
in a web browser. The main caveat is that this is a 39-image study; the results
are encouraging but need confirmation on larger, multi-hospital datasets.
\end{abstract}

\begin{IEEEkeywords}
medical image compression, lossless compression, entropy coding, asymmetric
numeral systems (ANS), finite state entropy (FSE), DICOM, high-bit-depth
imaging.
\end{IEEEkeywords}

% ============================================================
\section{Why Medical Images Need a Different Codec}
\label{sec:intro}

\IEEEPARstart{A}{modern} scanner produces a great deal of data. A single
digital breast tomosynthesis view can comprise 60 or more image slices and over
600\,MB; a full screening study can exceed 2\,GB uncompressed. Hospitals store
and move very large volumes of this data every day, so effective lossless
compression directly affects storage cost, network load, and how quickly images
appear in a viewer.

DICOM~\cite{dicom}, the medical imaging standard, already supports several
lossless formats---JPEG~2000~\cite{taubman2002}, HTJ2K~\cite{htj2k2019},
JPEG-LS~\cite{loco1}, and a basic run-length scheme. Each makes a different
trade-off between how small the file gets, how fast it decodes, and how complex
it is to implement~\cite{liu2017,clunie2000}. JPEG~2000 and HTJ2K compress well
but rely on fairly heavy machinery. JPEG-LS usually gets the smallest files.
Plain run-length coding is simple but barely compresses.

\subsection{The core problem: most coders assume 8-bit data}

This is the central observation behind the paper. The widely used entropy
coders, including Huffman, arithmetic coding, the LZ77 family, and the FSE/ANS
coder inside Zstandard~\cite{rfc8878}, were all designed for \textbf{8-bit byte
streams}, where there are only 256 possible symbols. Medical images do not fit
that assumption: a 16-bit pixel has up to 65,536 possible values.

When an 8-bit coder is applied to 16-bit data, it has two poor options:

\begin{enumerate}
  \item \textbf{Split each pixel into a high byte and a low byte} and code them
  separately. This doubles the number of coding steps and, more importantly,
  discards the relationship between the two bytes.
  \item \textbf{Cap the alphabet at 256 values} and store anything rarer as raw
  data, which wastes space.
\end{enumerate}

Why does splitting hurt? Consider a tiny prediction error like $-2$, $-1$, 0,
or $+1$. The low byte carries the real information, but the high byte is
almost always the same (just a sign-extension). A byte-by-byte coder still
spends bits encoding that high byte even though it is nearly fully
determined by the low one. On the test images, the information shared
between the two bytes ranges from 0.3 to 1.1 bits per pixel, which is the
ceiling on how much a byte-splitting coder is wasting, and MIC's measured
5--22\% advantage falls within that ceiling.

\paragraph{The residual stream} Throughout, a \emph{residual} is the
difference between a pixel and a prediction of it. Example: if the predictor
guesses 1020 and the true pixel is 1023, the residual is $+3$. Smooth regions
of an image produce many residuals near zero, and the same small values recur
frequently, which is exactly the pattern run-length and entropy coding exploit
well.

\subsection{What MIC is}

MIC keeps the residuals at their native 16-bit width and combines (i)~a simple
spatial predictor, (ii)~a 16-bit run-length encoder, and (iii)~an entropy coder
(FSE) extended to handle large alphabets. Plus a multi-state decoder that
improves decompression speed by letting the CPU work on several symbols at
once.

\subsection{Contributions, briefly}

\begin{itemize}
  \item A \textbf{16-bit-native pipeline} (Delta + 16-bit RLE + extended FSE)
  that outperforms Delta + Zstandard on the large majority of images (34 of 39,
  by about 10\% in geometric mean).
  \item An \textbf{entropy coder for large alphabets}---up to 65,535 symbols,
  versus the 4,096-symbol limit in Zstandard's FSE. Tables shrink
  automatically to fit the data (as small as ${\sim}2$\,KB for 8-bit content).
  \item A \textbf{multi-state decoder} that runs 2/4/8 independent decode chains
  to use idle CPU capacity, with no cost to compression ratio.
  \item An \textbf{evaluation} against HTJ2K, JPEG-LS, and Delta + Zstandard.
  \item A \textbf{20\,KB JavaScript decoder} plus a WebAssembly build, for
  client-side decoding in browsers.
\end{itemize}

\paragraph{Scope} This work covers lossless, single-frame, grayscale images
only. Color and multi-frame extensions, lossy modes, and progressive decoding
are out of scope.

% ============================================================
\section{Related Work}
\label{sec:related}

\textbf{The DICOM lossless codecs.} JPEG~2000~\cite{taubman2002} uses a wavelet
transform plus a sophisticated block coder (EBCOT): strong ratios, but slow
decoding. HTJ2K~\cite{htj2k2019,taubman2022} replaces the block coder with a
faster one and decodes considerably quicker. JPEG-LS~\cite{loco1} uses a small
predictor (Median Edge Detector) and Golomb-Rice coding; it is a widely
deployed standard and a strong ratio baseline. The basic DICOM run-length
scheme compresses little because it works on raw bytes with no prediction step.
Research codecs such as CALIC~\cite{wu1997} and FLIF~\cite{sneyers2016}
compress well on natural images but are not DICOM standards and lack browser
decoders.

These are the codecs MIC is measured against. The goal is not to outperform
them on every metric, but to explore a \textbf{simpler, residual-domain design}
aimed at high-bit-depth images, with strong speed and easy deployment.

\textbf{Entropy coding background.} Asymmetric Numeral Systems
(ANS)~\cite{duda2009,duda2013} is a family of entropy coders that replaces
arithmetic coding's interval math with simple integer state updates. FSE is a
fast, table-driven version of ANS~\cite{fse} made popular by
Zstandard~\cite{rfc8878}. The limitation is that these implementations are
tuned for byte streams (Zstandard caps FSE at 4,096 symbols).

\textbf{Parallel ANS.} Prior work has shown that ANS decoders run faster when
several independent streams are decoded at once (Giesen on interleaved
coders~\cite{giesen2014}; later work scaling this to many
streams~\cite{lin2023} and to GPUs~\cite{weissenberger2019}). MIC does not claim
to invent ANS or interleaving; it applies these known ideas inside a
16-bit-native medical codec with a specific large-alphabet design.

Table~\ref{tab:positioning} places MIC against the prior work across five
dimensions relevant to medical image compression.

\begin{table}[htbp]
\centering
\caption{Feature comparison of MIC with related codecs and entropy coding
methods. ``Entropy coder sees 16-bit residual as one symbol'' is a property
of the entropy stage, not of the codec's overall bit-depth support: JPEG-LS
and HTJ2K both handle 16-bit medical images, but their entropy stages
(Golomb--Rice on context bins for JPEG-LS, bit-plane coding for HTJ2K) do
not treat the predicted residual as a single 16-bit symbol. Source-line
counts are order-of-magnitude context only; the reference libraries also
implement standard profiles outside MIC's current scope.}
\label{tab:positioning}
\scriptsize
\setlength{\tabcolsep}{3pt}
\resizebox{\columnwidth}{!}{%
\begin{tabular}{p{2.7cm}ccccc}
\toprule
Feature & JPEG-LS & HTJ2K & Zstd & rANS & MIC \\
\midrule
Lossless medical & Yes & Yes & General & General & Yes \\
Entropy coder sees 16-bit residual as one symbol & No & No & No & Configurable & Yes \\
Multi-state ILP & No & No & No & Yes & Yes \\
Browser decoder & Via WASM port & Via WASM port & Partial & No & Purpose-built JS \\
Source lines & 5k & 20k & N/A & N/A & 1.1k \\
\bottomrule
\end{tabular}}
\end{table}

% ============================================================
\section{How MIC Works}
\label{sec:pipeline}

The pipeline is three stages, each simple enough to decode in one
straight-line pass with very few branches: raw 16-bit pixels go through
spatial prediction (subtract a guess based on neighbors), then 16-bit
run-length encoding (collapse repeated values), then FSE entropy coding
(remove the remaining redundancy), producing compressed bytes.

\subsection{Spatial prediction}

For each pixel, MIC predicts its value as the \textbf{average of the pixel
above and the pixel to the left}, then stores only the difference (the
residual). Pixels on the top edge or left edge use whatever single neighbor is
available; the very first pixel is stored as-is. Division rounds down, and
the decoder repeats the exact same arithmetic so reconstruction is exact.

This predictor was chosen because it is inexpensive, has almost no branches on
the decode side, and works well on medical images.

\paragraph{Other predictors we evaluated} We compared four predictors
(left-only, average, Paeth from PNG~\cite{png}, and MED from JPEG-LS~\cite{loco1})
through the full pipeline. MED gave the best ratio (about 2.4\% better than
average) but decoded 1.5--2$\times$ slower because of its three-way branch and a
diagonal dependency that blocks the fast branch-free loop. The average
predictor was preferred for its consistency and speed, so it is the default.
Table~\ref{tab:predictor-summary} summarizes the comparison.

\begin{table}[htbp]
\centering
\caption{Predictor ablation summary: geometric means and win counts across all
39 images.}
\label{tab:predictor-summary}
\begin{tabular}{lrrr}
\toprule
Predictor & Geo.\ Mean & Wins & Avg$\to$X \\
\midrule
Left-only & 3.22$\times$ & 2/39 & $-$4.4\% \\
Average (MIC) & 3.36$\times$ & 11/39 & baseline \\
Paeth & 3.39$\times$ & 0/39 & $+$0.8\% \\
MED (JPEG-LS) & 3.44$\times$ & 26/39 & $+$2.4\% \\
\bottomrule
\end{tabular}
\end{table}

\subsection{Handling large jumps: overflow coding}

Most residuals are small, but occasionally a pixel jumps far from its
prediction. Rather than widen the whole stream to 32 bits to handle rare large
values, MIC reserves one special code as a \textbf{delimiter}: small, in-range
residuals are stored directly as 16-bit codes (with zero mapped to a fixed
midpoint value); a large, out-of-range residual is stored as the delimiter
followed by the raw pixel value.

On decode, the reader checks each value: if it is the delimiter, the next
value is a raw pixel; otherwise it is a normal residual. The ranges are
designed not to overlap, so there is never any ambiguity. In normal 8--16-bit
clinical data the overflow path is rare; it exists to guarantee correctness,
not for the common case. Table~\ref{tab:overflow-map} shows the mapping for a
representative bit depth.

\begin{table}[htbp]
\centering
\caption{Overflow coding example for $d=12$ bits
(\texttt{deltaThreshold}$=2047$, delimiter $D=4095$).}
\label{tab:overflow-map}
\scriptsize
\begin{tabular}{lll}
\toprule
Residual $r$ & Emitted code(s) & Notes \\
\midrule
$-2$ & $2045$ & in-range \\
$-1$ & $2046$ & in-range \\
$0$  & $2047$ & in-range (= threshold) \\
$+1$ & $2048$ & in-range \\
$+2$ & $2049$ & in-range \\
$\pm 2047$ & $4095$, raw pixel & overflow \\
\bottomrule
\end{tabular}
\end{table}

\subsection{16-bit run-length encoding (RLE)}

The residual stream is turned into two kinds of runs: \textbf{same runs} (a
count plus one repeated value, used when a value appears at least three times
in a row) and \textbf{literal runs} (a count followed by that many values
copied verbatim; the values need not be distinct). A single header number tells
the two apart using a midpoint trick: small header values mean ``same run of
this length,'' large ones mean ``literal run.'' The longest single run is
${\sim}32{,}767$ symbols; longer stretches are split across several headers
automatically.

\textbf{Why 16-bit RLE matters.} Suppose 1,000 identical 16-bit residuals
appear in a row. MIC encodes that as just two numbers: a count and the value.
A \emph{byte}-oriented RLE sees the high and low bytes interleaved, so no
single byte repeats 1,000 times, so it is forced into two interleaved runs or
falls back to literals. Working at 16 bits captures the repetition that
byte-splitting hides.

\textbf{Fast decode.} Same runs (often 80\% of all runs on smooth images) are
fast-pathed: the decoder just hands back the cached value and decrements a
counter, without touching the compressed data at all.

\subsection{The entropy coder (FSE) for large alphabets}

\textbf{FSE in one paragraph.} FSE keeps a single integer ``state.'' To decode
one symbol, it looks up the current state in a table; the entry tells it which
symbol to output, how many bits to read next, and how to compute the next
state. So each symbol costs \textbf{one table lookup, one bit read, and one
addition}---no multiplies or divides. (Encoding runs the symbols in reverse
and builds the bitstream backwards; decoding reads forwards. That reversal is
inherent to how ANS works.)

MIC extends FSE to handle up to \textbf{65,535 symbols} instead of Zstandard's
4,096, because high-bit-depth residuals genuinely require a large alphabet. It
also \textbf{sizes its tables to the data}: an 8-bit image stored in 16-bit
containers (common for MR) shrinks the working tables from ${\sim}512$\,KB down
to about 2\,KB, which keeps everything in fast cache.

\textbf{Picking the table size automatically.} FSE has a ``tableLog'' knob that
sets how large and precise the probability table is. MIC selects it based on how
many distinct values it sees and how much data backs each value: it uses a
larger table only when there are many active symbols \emph{and} enough data to
estimate their probabilities reliably (for example, mammography images after
prediction). On the test set this adaptive choice correctly turned on the
larger table for exactly the nine images that benefit (1--10\% better ratio),
and decode speed is unaffected by the choice.
Table~\ref{tab:tablelog-ablation} shows the ablation across selected images.

\begin{table}[htbp]
\centering
\caption{TableLog ablation: forced TL=11/12/13 vs.\ adaptive selection
(selected images, ARM64 reference platform).}
\label{tab:tablelog-ablation}
\scriptsize
\begin{tabular}{lrrrr}
\toprule
Image & TL=11 & TL=12 & TL=13 & Adaptive \\
\midrule
CR  & 3.47$\times$ & 3.63$\times$ & 3.69$\times$ & 3.69$\times$ \\
MG1 & 8.00$\times$ & 8.57$\times$ & 8.79$\times$ & 8.79$\times$ \\
MG2 & 7.98$\times$ & 8.55$\times$ & 8.77$\times$ & 8.77$\times$ \\
MR2 & 3.14$\times$ & 3.24$\times$ & 3.28$\times$ & 3.28$\times$ \\
RG2 & 3.85$\times$ & 4.12$\times$ & 4.23$\times$ & 4.23$\times$ \\
RG3 & 5.64$\times$ & 5.95$\times$ & 6.08$\times$ & 6.08$\times$ \\
\bottomrule
\end{tabular}
\end{table}

% ============================================================
\section{Faster Decoding: Multi-State Decoding}
\label{sec:multistate}

\subsection{The bottleneck}

A standard single-state ANS decoder has an inherent limitation: each symbol
depends on the one before it. To decode symbol $N$ the decoder needs the state
produced by symbol $N-1$, so the work cannot start early. With a table lookup
taking roughly 4--5 CPU cycles, the loop is limited to about one symbol every
few cycles, even though modern CPUs have 4--6 execution ports that remain
largely idle. A single decode chain uses only one of them, leaving
${\sim}75$--83\% of the CPU's capacity unused.

\subsection{The approach}

MIC runs \textbf{several independent decoders at once} on interleaved
positions. With 8 chains, chain 0 handles positions $0, 8, 16, \ldots$; chain 1
handles $1, 9, 17, \ldots$; and so on. The chains share nothing during
decoding, each has its own state and its own stream of lookups, so the CPU's
out-of-order engine can keep all of them in flight simultaneously.

On the encoder side, the symbols are split into $N$ interleaved groups, each
compressed independently, and the resulting streams are concatenated. They
all share one decode table, and each chain's starting state is recorded in
the header. The decoder just writes each chain's output back to its
positions.

\subsection{Measured speedup}

In theory, $N$ chains could be up to $N\times$ faster. In practice the speedup
is smaller because the chains share some bookkeeping (refilling the bit
reader), the unrolled loop puts pressure on the instruction cache, and reading
the compressed data has its own bandwidth limit~\cite{williams2009}. Measured
on the FSE stage alone, \textbf{4 states} give $+$36\% on ARM64 and $+$12\% on
AMD64 (relative to single state), and \textbf{8 states} a further $+$5\% on
ARM64 and $+$1\% on AMD64, with \textbf{no change to the compressed file}.
ARM64 continues to improve as chains double; AMD64 flattens out earlier because
its core already extracts substantial parallelism from a single chain.
Table~\ref{tab:latency} shows the simple latency model against the measured
ARM64 speedup.

\begin{table}[htbp]
\centering
\caption{Multi-state decoder latency model and empirical speedup.}
\label{tab:latency}
\begin{tabular}{clrr}
\toprule
States ($k$) & Amortized latency & Theoretical & Empirical (ARM64) \\
\midrule
1 & $L$ cy/sym & 1.0$\times$ & 1.0$\times$ \\
2 & $L/2$ cy/sym & 2.0$\times$ & 1.3$\times$ \\
4 & $L/4$ cy/sym & 4.0$\times$ & 1.5$\times$ \\
8 & $L/8$ cy/sym & 8.0$\times$ & 1.6$\times$ \\
\bottomrule
\end{tabular}
\end{table}

The decoder signals its mode with a short magic header (a reserved byte
value that cannot be a valid single-state setting), so single-, two-, four-,
and eight-state files all coexist without confusion. On AMD64 and ARM64 the
inner loop uses each architecture's native variable-shift instructions; other
platforms get a pure-Go fallback. Table~\ref{tab:format} summarises the
on-disk layout of a single-frame stream.

\begin{table}[htbp]
\centering
\caption{Single-frame MIC bitstream layout. All multi-byte fields are
little-endian. \emph{States} $k$ is the number of interleaved FSE chains
($k \in \{1,2,4,8\}$).}
\label{tab:format}
\scriptsize
\begin{tabular}{p{2.4cm}rp{4.4cm}}
\toprule
Field & Size & Meaning \\
\midrule
Magic / mode byte & 1 B & Single-state: encodes \texttt{tableLog} (value 5--17). Multi-state: \texttt{0xFF}. \\
Multi-state tag & 1 B & For multi-state streams, one of \texttt{0x02}, \texttt{0x04}, \texttt{0x84} indicating $k=2$, 4, or 8. \\
Symbol count & 4 B & Number of decoded RLE symbols. \\
TableLog & 1 B & Decode-table size: $2^{\textit{tableLog}}$ entries (multi-state path only; the single-state path overloads byte~0). \\
Normalized counts & variable & FSE normalised symbol frequencies, serialised with the standard variable-length encoding. \\
Per-chain initial states & $k \times 2$ B & Starting state for each FSE chain. \\
Per-chain stream lengths & $k \times 4$ B & Compressed-stream length per chain (omitted when $k=1$). \\
Per-chain compressed streams & variable & Concatenated chain payloads. \\
\bottomrule
\end{tabular}
\end{table}

% ============================================================
\section{Evaluation Methodology}
\label{sec:methodology}

\textbf{Dataset.} 39 de-identified DICOM images across eleven modalities (CT,
MR, CR, X-ray, mammography, digital radiography, nuclear medicine, PET,
fluoroscopy, secondary capture, and tomosynthesis), drawn from the public NEMA
WG-04 sample sets~\cite{nema_wg04}, the NEMA 1997 CD, a public
breast-tomosynthesis case~\cite{clunie_tomo}, and additional grayscale slices
from public GDCM and TCIA collections. It is a varied set but small relative to
a real clinical archive, so the results are reported as ``on this 39-image
dataset'' rather than as general claims. Table~\ref{tab:dataset} lists the
images.

\begin{table}[htbp]
\centering
\caption{Test dataset: 39 clinical DICOM images spanning eleven modalities.}
\label{tab:dataset}
\scriptsize
\begin{tabular}{llrr}
\toprule
Image & Modality & Dimensions & Raw (MB) \\
\midrule
MR   & Brain MRI      & $256{\times}256$   & 0.13 \\
CT   & CT scan        & $512{\times}512$   & 0.50 \\
CR   & Comp.\ radiography & $2140{\times}1760$ & 7.18 \\
XR   & X-ray          & $2048{\times}2577$ & 10.1 \\
MG1  & Mammography    & $2457{\times}1996$ & 9.35 \\
MG2  & Mammography    & $2457{\times}1996$ & 9.35 \\
MG3  & Mammography    & $4774{\times}3064$ & 27.3 \\
MG4  & Mammography    & $4096{\times}3328$ & 26.0 \\
CT1  & CT scan        & $512{\times}512$   & 0.50 \\
CT2  & CT scan        & $512{\times}512$   & 0.50 \\
MG-N & Mammography    & $3064{\times}4664$ & 27.3 \\
MR1  & Brain MRI      & $512{\times}512$   & 0.50 \\
MR2  & Brain MRI      & $1024{\times}1024$ & 2.00 \\
MR3  & Brain MRI      & $512{\times}512$   & 0.50 \\
MR4  & Brain MRI      & $512{\times}512$   & 0.50 \\
NM1  & Nuclear med.   & $256{\times}1024$  & 0.50 \\
RG1  & Radiography    & $1841{\times}1955$ & 6.86 \\
RG2  & Radiography    & $1760{\times}2140$ & 7.18 \\
RG3  & Radiography    & $1760{\times}1760$ & 5.91 \\
SC1  & Sec.~capture   & $2048{\times}2487$ & 9.71 \\
XA1  & Fluoroscopy    & $1024{\times}1024$ & 2.00 \\
CT\_BRAIN    & CT scan             & $512{\times}512$   & 0.50 \\
CT\_ANKLE    & CT scan             & $512{\times}512$   & 0.50 \\
CT\_ORT      & CT scan             & $512{\times}512$   & 0.50 \\
CT\_CHEST    & CT scan             & $400{\times}512$   & 0.39 \\
MR\_HEAD     & Brain MRI           & $256{\times}256$   & 0.13 \\
MR\_INTERA   & Brain MRI           & $1024{\times}1024$ & 2.00 \\
DX\_HAND     & Digital radiography & $1480{\times}1410$ & 3.98 \\
DX\_CHEST    & Digital radiography & $1416{\times}1420$ & 3.84 \\
CR\_THORAX   & Comp.\ radiography  & $2076{\times}1876$ & 7.43 \\
PET1         & PET                 & $256{\times}256$   & 0.13 \\
PET2         & PET                 & $256{\times}256$   & 0.13 \\
CT\_LUNG     & CT scan             & $512{\times}512$   & 0.50 \\
CT\_PANCREAS & CT scan             & $512{\times}512$   & 0.50 \\
MR\_BRAIN    & Brain MRI           & $320{\times}256$   & 0.16 \\
MR\_BREAST   & Breast MRI          & $256{\times}256$   & 0.13 \\
MR\_PROSTATE & Prostate MRI        & $320{\times}320$   & 0.20 \\
PET\_PSMA    & PET                 & $200{\times}200$   & 0.08 \\
PET\_LUNG    & PET                 & $200{\times}200$   & 0.08 \\
\bottomrule
\end{tabular}
\end{table}

\textbf{Baselines.} HTJ2K (via OpenJPH~\cite{openjph}), JPEG-LS (via
CharLS~\cite{charls}), and Delta + Zstandard~\cite{rfc8878} at its maximum
level. All were called in-process (no file or subprocess overhead) for a fair
speed comparison.

\textbf{Metrics.} Compression ratio (uncompressed $\div$ compressed), and
decode and encode throughput in MB/s of uncompressed pixels.

\textbf{Platforms.} Two reference machines: an ARM64 (Apple M4 Pro) and an
AMD64 (AWS Intel Xeon 6). Go code used the standard compiler; C components were
built with \texttt{-O3}. Each benchmark ran the full image 10 times;
run-to-run variance stayed under 2\% on ARM64 and 5\% on AMD64.

The ``MIC'' column in the result tables is the fastest available C decoder on
each platform: the eight-state decoder, with AVX-512 acceleration of the
run-length and prediction-reversal stages on AMD64. Both read identical
compressed bytes; only the vector width differs.

\textbf{Browser decoder.} A 20\,KB pure-JavaScript ES module (zero npm
dependencies) plus a 2.5\,MB Go WebAssembly build. The JS decoder
auto-detects 1/2/4/8-state streams and was benchmarked under Node.js on the
ARM64 machine.

% ============================================================
\section{Results}
\label{sec:results}

\subsection{Compression ratio}

JPEG-LS produces the smallest files on nearly every image (37 of 39); its
context-adaptive predictor is difficult to beat on pure ratio. MIC, however, is
not optimizing ratio alone; it targets the best balance of ratio, speed, and
simplicity. Across the dataset MIC reaches a \textbf{3.36}$\times$
geometric-mean ratio: about 88\% of JPEG-LS (3.84$\times$) and within about 2\%
of HTJ2K (3.42$\times$), while decoding faster than both. Mammography compresses
best (up to 8.79$\times$) because its smooth tissue produces long near-zero
runs, and studies with large uniform backgrounds reach the highest single-image
ratios (PET\_PSMA at 10.12$\times$, CT\_ANKLE at 9.48$\times$).
Table~\ref{tab:ratios} gives the per-image ratios for all variants.

\begin{table*}[t]
\centering
\caption{Lossless compression ratios: all codec variants, 39 images. Bold =
best ratio per row.}
\label{tab:ratios}
\scriptsize
\begin{tabular}{lrrrrrrrr}
\toprule
Image & Raw (MB) & $\Delta$+Zstd-19 & MIC & Wavelet & PICS-4 & PICS-8 & HTJ2K & JPEG-LS \\
\midrule
MR             & 0.13 & 1.95$\times$ & 2.35$\times$ & 2.38$\times$ & 2.28$\times$ & 2.21$\times$ & 2.38$\times$ & \textbf{2.52$\times$} \\
CT             & 0.50 & 2.05$\times$ & 2.24$\times$ & 1.67$\times$ & 2.15$\times$ & 1.96$\times$ & 1.77$\times$ & \textbf{2.68$\times$} \\
CR             & 7.18 & 3.24$\times$ & 3.69$\times$ & 3.81$\times$ & 3.70$\times$ & 3.71$\times$ & 3.77$\times$ & \textbf{3.96$\times$} \\
XR             & 10.1 & 1.43$\times$ & 1.74$\times$ & \textbf{1.76$\times$} & 1.75$\times$ & 1.75$\times$ & 1.67$\times$ & \textbf{1.76$\times$} \\
MG1            & 9.35 & 7.20$\times$ & 8.79$\times$ & 8.66$\times$ & 8.84$\times$ & 8.87$\times$ & 8.25$\times$ & \textbf{8.91$\times$} \\
MG2            & 9.35 & 7.21$\times$ & 8.77$\times$ & 8.65$\times$ & 8.83$\times$ & 8.85$\times$ & 8.24$\times$ & \textbf{8.90$\times$} \\
MG3            & 27.3 & 2.09$\times$ & 2.24$\times$ & 2.31$\times$ & 2.31$\times$ & 2.33$\times$ & 2.22$\times$ & \textbf{2.38$\times$} \\
MG4            & 26.0 & 3.10$\times$ & 3.47$\times$ & 3.59$\times$ & 3.58$\times$ & 3.62$\times$ & 3.51$\times$ & \textbf{3.71$\times$} \\
CT1            & 0.50 & 2.47$\times$ & 2.79$\times$ & 2.48$\times$ & 2.54$\times$ & 2.29$\times$ & 2.70$\times$ & \textbf{3.19$\times$} \\
CT2            & 0.50 & 3.24$\times$ & 3.48$\times$ & 2.87$\times$ & 3.11$\times$ & 2.72$\times$ & 3.29$\times$ & \textbf{4.54$\times$} \\
MG-N           & 27.3 & 2.04$\times$ & 2.24$\times$ & 2.32$\times$ & 2.31$\times$ & 2.34$\times$ & 2.23$\times$ & \textbf{2.38$\times$} \\
MR1            & 0.50 & 1.79$\times$ & 2.09$\times$ & 2.14$\times$ & 2.10$\times$ & 2.08$\times$ & 2.12$\times$ & \textbf{2.30$\times$} \\
MR2            & 2.00 & 2.77$\times$ & 3.28$\times$ & 3.34$\times$ & 3.31$\times$ & 3.31$\times$ & 3.35$\times$ & \textbf{3.52$\times$} \\
MR3            & 0.50 & 3.47$\times$ & 3.92$\times$ & 4.09$\times$ & 3.89$\times$ & 3.83$\times$ & 4.33$\times$ & \textbf{4.51$\times$} \\
MR4            & 0.50 & 3.58$\times$ & 4.12$\times$ & 4.18$\times$ & 4.09$\times$ & 4.03$\times$ & 4.21$\times$ & \textbf{4.49$\times$} \\
NM1            & 0.50 & 4.88$\times$ & 5.15$\times$ & 5.01$\times$ & 5.25$\times$ & 5.28$\times$ & 5.76$\times$ & \textbf{6.28$\times$} \\
RG1            & 6.86 & 1.45$\times$ & 1.70$\times$ & 1.70$\times$ & 1.70$\times$ & 1.69$\times$ & 1.63$\times$ & \textbf{1.72$\times$} \\
RG2            & 7.18 & 3.69$\times$ & 4.23$\times$ & 4.32$\times$ & 4.28$\times$ & 4.30$\times$ & 4.32$\times$ & \textbf{4.51$\times$} \\
RG3            & 5.91 & 5.46$\times$ & 6.08$\times$ & 6.82$\times$ & 6.11$\times$ & 6.12$\times$ & 6.99$\times$ & \textbf{7.31$\times$} \\
SC1            & 9.71 & 3.46$\times$ & 3.71$\times$ & 3.70$\times$ & 3.73$\times$ & 3.73$\times$ & 3.85$\times$ & \textbf{4.73$\times$} \\
XA1            & 2.00 & 4.37$\times$ & 5.01$\times$ & 4.94$\times$ & 5.04$\times$ & 5.03$\times$ & 4.88$\times$ & \textbf{5.39$\times$} \\
CT\_BRAIN      & 0.50 & 3.08$\times$ & 3.06$\times$ & 2.89$\times$ & 2.77$\times$ & 2.47$\times$ & 3.39$\times$ & \textbf{4.28$\times$} \\
CT\_ANKLE      & 0.50 & 9.30$\times$ & \textbf{9.48$\times$} & 5.02$\times$ & 9.30$\times$ & 9.14$\times$ & 4.97$\times$ & 6.87$\times$ \\
CT\_ORT        & 0.50 & 2.36$\times$ & 2.68$\times$ & 2.38$\times$ & 2.51$\times$ & 2.26$\times$ & 2.56$\times$ & \textbf{3.00$\times$} \\
CT\_CHEST      & 0.39 & 2.46$\times$ & 2.82$\times$ & 2.10$\times$ & 2.52$\times$ & 2.21$\times$ & 2.23$\times$ & \textbf{3.22$\times$} \\
MR\_HEAD       & 0.13 & 2.11$\times$ & 2.61$\times$ & 2.65$\times$ & 2.57$\times$ & 2.53$\times$ & 2.67$\times$ & \textbf{2.86$\times$} \\
MR\_INTERA     & 2.00 & 4.71$\times$ & 3.68$\times$ & 4.05$\times$ & 3.75$\times$ & 3.78$\times$ & 5.04$\times$ & \textbf{5.08$\times$} \\
DX\_HAND       & 3.98 & 1.99$\times$ & 2.24$\times$ & 2.35$\times$ & 2.26$\times$ & 2.25$\times$ & 2.37$\times$ & \textbf{2.48$\times$} \\
DX\_CHEST      & 3.84 & 1.43$\times$ & 1.66$\times$ & \textbf{1.82$\times$} & 1.65$\times$ & 1.63$\times$ & 1.63$\times$ & 1.70$\times$ \\
CR\_THORAX     & 7.43 & 1.61$\times$ & 1.82$\times$ & 1.81$\times$ & 1.84$\times$ & 1.83$\times$ & 1.81$\times$ & \textbf{1.90$\times$} \\
PET1           & 0.13 & 2.37$\times$ & 2.74$\times$ & 2.81$\times$ & 2.64$\times$ & 2.52$\times$ & 3.02$\times$ & \textbf{3.21$\times$} \\
PET2           & 0.13 & 3.39$\times$ & 3.39$\times$ & 3.48$\times$ & 3.16$\times$ & 2.58$\times$ & 4.37$\times$ & \textbf{4.74$\times$} \\
CT\_LUNG       & 0.50 & 2.40$\times$ & 2.73$\times$ & 2.45$\times$ & 2.50$\times$ & 2.25$\times$ & 2.67$\times$ & \textbf{3.15$\times$} \\
CT\_PANCREAS   & 0.50 & 2.09$\times$ & 2.36$\times$ & 1.76$\times$ & 2.18$\times$ & 1.98$\times$ & 1.85$\times$ & \textbf{2.70$\times$} \\
MR\_BRAIN      & 0.16 & 6.17$\times$ & 7.27$\times$ & 6.75$\times$ & 7.31$\times$ & 7.24$\times$ & 7.62$\times$ & \textbf{8.46$\times$} \\
MR\_BREAST     & 0.13 & 3.31$\times$ & 4.01$\times$ & 3.94$\times$ & 4.00$\times$ & 3.91$\times$ & 4.10$\times$ & \textbf{4.48$\times$} \\
MR\_PROSTATE   & 0.20 & 2.04$\times$ & 2.30$\times$ & 2.34$\times$ & 2.28$\times$ & 2.26$\times$ & 2.49$\times$ & \textbf{2.65$\times$} \\
PET\_PSMA      & 0.08 & 11.88$\times$ & 10.12$\times$ & 7.67$\times$ & 1.20$\times$ & 1.18$\times$ & 13.91$\times$ & \textbf{15.78$\times$} \\
PET\_LUNG      & 0.08 & 4.56$\times$ & 2.97$\times$ & 2.97$\times$ & 1.00$\times$ & 1.03$\times$ & 5.21$\times$ & \textbf{5.62$\times$} \\
\midrule
\textbf{Geomean} & --- & \textbf{3.06$\times$} & \textbf{3.36$\times$} & \textbf{3.21$\times$} & \textbf{3.04$\times$} & \textbf{2.95$\times$} & \textbf{3.42$\times$} & \textbf{3.84$\times$} \\
\bottomrule
\end{tabular}
\end{table*}

\subsection{Encoding speed}

MIC encodes faster than the others on most images, on both platforms,
because its encoder is a single pass (predict, run-length, table-driven
entropy code) while HTJ2K and JPEG-LS do more work per pixel. On
\textbf{ARM64}, MIC is fastest on 38 of 39 images: roughly 1.0--3.3$\times$
HTJ2K, 1.0--7.1$\times$ JPEG-LS, and 30--320$\times$ Delta + Zstandard (whose
max-level LZ77 search accounts for its low throughput); HTJ2K edges ahead only
on the tiny PET\_LUNG study. On \textbf{AMD64}, MIC is fastest on 32 of 39;
HTJ2K leads on the two largest mammography images, one radiography image, and
four of the smaller new-corpus studies, because the AMD64 baselines are
themselves heavily vectorized and per-image table-build overhead weighs more on
tiny images.
Table~\ref{tab:enc} gives the per-image encoding throughput.

\begin{table*}[htbp]
\centering
\caption{Single-threaded encoding throughput (MB/s). MIC is the
eight-state C encoder. Bold marks the fastest codec per row on each
platform.}
\label{tab:enc}
\scriptsize
\begin{tabular}{l|rrrr|rrrr}
\toprule
 & \multicolumn{4}{c|}{ARM64} & \multicolumn{4}{c}{AMD64} \\
\cmidrule(lr){2-5}\cmidrule(lr){6-9}
Image & MIC & $\Delta$+Zstd & HTJ2K & JPEG-LS & MIC & $\Delta$+Zstd & HTJ2K & JPEG-LS \\
\midrule
MR             & \textbf{457} & 7 & 236 & 91 & \textbf{344} & 5 & 246 & 62 \\
CT             & \textbf{476} & 7 & 183 & 126 & \textbf{280} & 5 & 207 & 97 \\
CR             & \textbf{525} & 2 & 186 & 110 & \textbf{322} & 2 & 286 & 75 \\
XR             & \textbf{661} & 3 & 229 & 113 & \textbf{400} & 4 & 256 & 84 \\
MG1            & \textbf{1{,}057} & 11 & 414 & 277 & 532 & 8 & \textbf{632} & 211 \\
MG2            & \textbf{1{,}047} & 12 & 407 & 284 & 552 & 8 & \textbf{629} & 209 \\
MG3            & \textbf{639} & 2 & 211 & 132 & \textbf{352} & 2 & 245 & 104 \\
MG4            & \textbf{883} & 5 & 342 & 196 & \textbf{457} & 6 & 408 & 151 \\
CT1            & \textbf{615} & 9 & 243 & 176 & \textbf{321} & 7 & 271 & 125 \\
CT2            & \textbf{536} & 7 & 204 & 172 & \textbf{285} & 5 & 271 & 124 \\
MG-N           & \textbf{639} & 2 & 214 & 140 & \textbf{358} & 2 & 248 & 104 \\
MR1            & \textbf{616} & 9 & 205 & 114 & \textbf{367} & 6 & 241 & 86 \\
MR2            & \textbf{685} & 7 & 216 & 118 & \textbf{409} & 4 & 309 & 85 \\
MR3            & \textbf{768} & 19 & 264 & 167 & \textbf{441} & 12 & 354 & 123 \\
MR4            & \textbf{585} & 7 & 212 & 163 & \textbf{367} & 5 & 312 & 142 \\
NM1            & \textbf{633} & 6 & 208 & 162 & \textbf{340} & 5 & 313 & 115 \\
RG1            & \textbf{595} & 4 & 218 & 84 & \textbf{396} & 5 & 252 & 66 \\
RG2            & \textbf{737} & 4 & 263 & 151 & \textbf{425} & 3 & 360 & 100 \\
RG3            & \textbf{757} & 6 & 276 & 183 & 389 & 4 & \textbf{412} & 126 \\
SC1            & \textbf{667} & 5 & 225 & 203 & \textbf{416} & 4 & 306 & 169 \\
XA1            & \textbf{612} & 5 & 203 & 133 & \textbf{356} & 4 & 303 & 98 \\
CT\_BRAIN      & \textbf{352} & 5 & 168 & 147 & 239 & 4 & \textbf{253} & 112 \\
CT\_ANKLE      & \textbf{955} & 15 & 290 & 327 & \textbf{450} & 11 & 358 & 261 \\
CT\_ORT        & \textbf{644} & 10 & 198 & 138 & \textbf{350} & 7 & 269 & 112 \\
CT\_CHEST      & \textbf{463} & 6 & 177 & 133 & \textbf{267} & 4 & 234 & 100 \\
MR\_HEAD       & \textbf{430} & 6 & 211 & 106 & \textbf{288} & 4 & 252 & 65 \\
MR\_INTERA     & \textbf{415} & 4 & 180 & 115 & \textbf{269} & 3 & 268 & 89 \\
DX\_HAND       & \textbf{587} & 4 & 218 & 102 & \textbf{378} & 3 & 261 & 65 \\
DX\_CHEST      & \textbf{636} & 6 & 209 & 109 & \textbf{389} & 6 & 241 & 77 \\
CR\_THORAX     & \textbf{585} & 4 & 226 & 91 & \textbf{397} & 4 & 261 & 68 \\
PET1           & \textbf{479} & 7 & 240 & 129 & \textbf{295} & 5 & 252 & 111 \\
PET2           & \textbf{550} & 6 & 236 & 145 & \textbf{269} & 4 & 256 & 103 \\
CT\_LUNG       & \textbf{586} & 10 & 195 & 141 & \textbf{337} & 7 & 258 & 110 \\
CT\_PANCREAS   & \textbf{454} & 7 & 193 & 130 & \textbf{292} & 5 & 240 & 98 \\
MR\_BRAIN      & \textbf{607} & 7 & 229 & 191 & 338 & 6 & \textbf{351} & 141 \\
MR\_BREAST     & \textbf{522} & 5 & 241 & 141 & \textbf{292} & 4 & 290 & 96 \\
MR\_PROSTATE   & \textbf{465} & 8 & 195 & 97 & \textbf{313} & 6 & 250 & 66 \\
PET\_PSMA      & \textbf{434} & 14 & 423 & 426 & 296 & 10 & \textbf{377} & 364 \\
PET\_LUNG      & 232 & 6 & \textbf{245} & 235 & 144 & 5 & \textbf{260} & 167 \\
\midrule
\textbf{Geomean} & \textbf{582} & \textbf{6} & \textbf{231} & \textbf{148} & \textbf{343} & \textbf{5} & \textbf{293} & \textbf{110} \\
\bottomrule
\end{tabular}
\end{table*}

\subsection{Decoding speed}

Decoding is where MIC performs best. Its decoder is dominated by the compact
FSE loop (one lookup, one bit read, one add) with no hard-to-predict branches.
HTJ2K and JPEG-LS both have data-dependent branches that cause CPU
mispredictions. Zstandard is the most comparable codec (same predictor, also
uses FSE) but works on bytes, so it pays the 2$\times$ byte-splitting cost on
deep data. On \textbf{ARM64}, MIC is fastest on all 39 images: roughly
1.2--2.1$\times$ HTJ2K, 1.0--2.0$\times$ Zstandard, and 1.1--5.1$\times$
JPEG-LS, with the largest margins on the smaller images. On \textbf{AMD64}, MIC
is fastest on 36 of 39 images; on the three exceptions (CT\_BRAIN, MR\_INTERA,
and the tiny PET\_LUNG study) in-process Zstandard is marginally faster. Across
its wins MIC runs about 1.0--1.8$\times$ HTJ2K and 1.9--7.1$\times$ JPEG-LS, and
stays within 0.7--1.8$\times$ of Zstandard; switching from four to eight states
(plus the AVX-512 widening of the post-entropy stages) is what turned several
remaining close cases into MIC wins.

MIC therefore delivers both \textbf{a better ratio than Zstandard on most
images} (34 of 39) and \textbf{faster decoding on nearly every image}, which
reflects the 16-bit-native design. JPEG-LS retains about a 14\% ratio
advantage, so the choice between them comes down to decode speed versus storage
cost. Per-image numbers are in Table~\ref{tab:decomp}, and the entropy stage in
isolation is in Table~\ref{tab:fse-combined}.

\begin{table*}[htbp]
\centering
\caption{Single-threaded decompression throughput (MB/s). MIC is
the eight-state C decoder (AVX-512 acceleration of the RLE and delta
inverse stages on AMD64). Bold marks the fastest codec per row on
each platform.}
\label{tab:decomp}
\scriptsize
\begin{tabular}{l|rrrr|rrrr}
\toprule
 & \multicolumn{4}{c|}{ARM64} & \multicolumn{4}{c}{AMD64} \\
\cmidrule(lr){2-5}\cmidrule(lr){6-9}
Image & MIC & $\Delta$+Zstd & HTJ2K & JPEG-LS & MIC & $\Delta$+Zstd & HTJ2K & JPEG-LS \\
\midrule
MR             & \textbf{649} & 389 & 328 & 137 & \textbf{614} & 360 & 453 & 91 \\
CT             & \textbf{495} & 406 & 299 & 158 & \textbf{444} & 393 & 329 & 110 \\
CR             & \textbf{580} & 527 & 316 & 181 & \textbf{570} & 479 & 445 & 124 \\
XR             & \textbf{668} & 512 & 324 & 131 & \textbf{626} & 509 & 393 & 88 \\
MG1            & \textbf{762} & 685 & 612 & 481 & \textbf{916} & 596 & 913 & 336 \\
MG2            & \textbf{753} & 692 & 622 & 492 & \textbf{942} & 599 & 902 & 334 \\
MG3            & \textbf{600} & 433 & 314 & 175 & \textbf{547} & 362 & 392 & 118 \\
MG4            & \textbf{788} & 504 & 508 & 246 & \textbf{686} & 469 & 634 & 158 \\
CT1            & \textbf{600} & 450 & 357 & 224 & \textbf{535} & 432 & 389 & 142 \\
CT2            & \textbf{570} & 471 & 342 & 212 & \textbf{522} & 458 & 383 & 131 \\
MG-N           & \textbf{600} & 408 & 289 & 180 & \textbf{549} & 357 & 398 & 118 \\
MR1            & \textbf{734} & 365 & 342 & 149 & \textbf{586} & 381 & 392 & 93 \\
MR2            & \textbf{700} & 445 & 363 & 211 & \textbf{682} & 460 & 486 & 131 \\
MR3            & \textbf{839} & 523 & 432 & 289 & \textbf{742} & 497 & 520 & 192 \\
MR4            & \textbf{692} & 466 & 326 & 220 & \textbf{664} & 496 & 433 & 152 \\
NM1            & \textbf{778} & 499 & 384 & 260 & \textbf{597} & 515 & 485 & 171 \\
RG1            & \textbf{469} & 383 & 318 & 123 & \textbf{547} & 340 & 363 & 80 \\
RG2            & \textbf{652} & 510 & 383 & 225 & \textbf{716} & 503 & 544 & 154 \\
RG3            & \textbf{685} & 642 & 432 & 294 & \textbf{719} & 564 & 641 & 195 \\
SC1            & \textbf{699} & 511 & 366 & 263 & \textbf{673} & 483 & 457 & 184 \\
XA1            & \textbf{662} & 501 & 346 & 244 & \textbf{670} & 541 & 501 & 165 \\
CT\_BRAIN      & \textbf{528} & 466 & 326 & 190 & 420 & \textbf{477} & 354 & 122 \\
CT\_ANKLE      & \textbf{853} & 587 & 452 & 426 & \textbf{774} & 562 & 537 & 282 \\
CT\_ORT        & \textbf{538} & 442 & 348 & 190 & \textbf{565} & 430 & 384 & 124 \\
CT\_CHEST      & \textbf{500} & 445 & 297 & 178 & \textbf{446} & 428 & 333 & 110 \\
MR\_HEAD       & \textbf{595} & 372 & 280 & 167 & \textbf{679} & 371 & 406 & 100 \\
MR\_INTERA     & \textbf{590} & 563 & 304 & 178 & 454 & \textbf{527} & 436 & 133 \\
DX\_HAND       & \textbf{500} & 414 & 316 & 147 & \textbf{570} & 416 & 362 & 101 \\
DX\_CHEST      & \textbf{476} & 352 & 307 & 121 & \textbf{547} & 349 & 358 & 83 \\
CR\_THORAX     & \textbf{552} & 403 & 346 & 129 & \textbf{561} & 365 & 365 & 84 \\
PET1           & \textbf{732} & 416 & 377 & 197 & \textbf{679} & 384 & 389 & 120 \\
PET2           & \textbf{656} & 515 & 381 & 227 & \textbf{487} & 438 & 389 & 141 \\
CT\_LUNG       & \textbf{522} & 316 & 348 & 189 & \textbf{534} & 420 & 383 & 122 \\
CT\_PANCREAS   & \textbf{517} & 451 & 316 & 172 & \textbf{439} & 390 & 337 & 110 \\
MR\_BRAIN      & \textbf{866} & 567 & 422 & 375 & \textbf{708} & 514 & 641 & 221 \\
MR\_BREAST     & \textbf{640} & 545 & 350 & 221 & \textbf{759} & 449 & 510 & 156 \\
MR\_PROSTATE   & \textbf{679} & 422 & 353 & 165 & \textbf{545} & 386 & 408 & 103 \\
PET\_PSMA      & \textbf{759} & 678 & 475 & 684 & \textbf{747} & 655 & 539 & 395 \\
PET\_LUNG      & \textbf{502} & 498 & 283 & 278 & 365 & \textbf{508} & 363 & 170 \\
\midrule
\textbf{Geomean} & \textbf{631} & \textbf{473} & \textbf{359} & \textbf{215} & \textbf{598} & \textbf{452} & \textbf{445} & \textbf{140} \\
\bottomrule
\end{tabular}
\end{table*}

\begin{table}[htbp]
\centering
\caption{FSE decompression throughput (MB/s) for representative modalities,
single-thread pure-Go decoder. Increasing the number of independent state
chains exposes more instruction-level parallelism to the out-of-order
engine. Per-image rows are ten representative modalities; both geometric
means are over all 39 images.}
\label{tab:fse-combined}
\scriptsize
\setlength{\tabcolsep}{4pt}
\begin{tabular}{lrrrrrr}
\toprule
 & \multicolumn{3}{c}{ARM64 (MB/s)} & \multicolumn{3}{c}{AMD64 (MB/s)} \\
\cmidrule(lr){2-4}\cmidrule(lr){5-7}
Image & 1-st & 4-st & 8-st & 1-st & 4-st & 8-st \\
\midrule
MR       & 1,066 & 1,576 & 1,068 &   656 & 1,007 &   847 \\
CT       & 1,001 & 1,549 & 2,142 &   972 & 1,034 & 1,362 \\
CR       & 2,941 & 4,882 & 6,245 & 1,456 & 1,844 & 1,709 \\
XR       & 3,597 & 6,081 & 6,030 & 1,745 & 2,341 & 2,014 \\
MG1      & 3,558 & 4,865 & 4,885 & 2,216 & 1,487 & 1,975 \\
MG-N     & 3,949 & 6,647 & 6,037 & 1,940 & 2,545 & 2,415 \\
NM1      & 1,971 & 2,661 & 2,570 &   917 & 1,346 & 1,356 \\
RG1      & 1,768 & 4,356 & 5,332 & 1,239 & 1,709 & 1,869 \\
DX\_HAND & 2,120 & 3,163 & 4,816 & 1,245 & 1,833 & 1,520 \\
PET1     & 1,105 & 1,252 & 1,530 &   575 &   763 &   616 \\
\midrule
\textbf{Geomean (39)} & \textbf{1,742} & \textbf{2,372} & \textbf{2,492} &
\textbf{1,093} & \textbf{1,228} & \textbf{1,246} \\
\bottomrule
\end{tabular}
\end{table}

\subsection{Other MIC variants (for reference)}

The main tables use one MIC column, but the codebase includes several variants
that exist for different deployment needs rather than to compete on raw speed.
The \textbf{pure Go} build runs the entire pipeline in Go with no C dependency,
at roughly 40--60\% of the C decoder's speed; this is useful where C cannot be
linked (mobile, WebAssembly, locked-down runtimes). The
\textbf{1/2/4-state decoders} are the same C decoder with fewer chains; they
serve as baselines for the multi-state comparison. \textbf{Wavelet-based MIC}
replaces the predictor with a reversible 5/3 wavelet transform~\cite{legall1988}
feeding the same back end; it is competitive on a few images but generally
slower on this corpus, since smooth medical images do not reward the wavelet's
additional cost. The eight-state C decoder is both the simplest and the fastest
configuration on this dataset.

\subsection{Multi-threaded decoding}

MIC can also split an image into horizontal strips and decode them on
separate threads. This scales nearly linearly up to four threads on both
platforms and keeps scaling to eight threads on images above ${\sim}1$\,MB,
hitting about \textbf{4.3\,GB/s} on the largest mammography images on ARM64.
Small images do not benefit: thread setup costs more than it saves, and on the
two 0.08\,MB PET images the strips fall back toward raw storage, so strip
parallelism does not apply. We report this separately because the other codecs
run single-threaded in our pipeline, so a direct comparison would not be fair.
Table~\ref{tab:pics} gives the per-image strip-parallel numbers.

\begin{table*}[htbp]
\centering
\caption{Multi-threaded decompression throughput (MB/s) for MIC's
strip-parallel mode. The ``MIC'' column reproduces the single-threaded
eight-state baseline from Table~\ref{tab:decomp} exactly; $N$-strip columns
decode the same compressed bytes in $N$ parallel threads. Both ARM64 and AMD64
columns cover all 39 images.}
\label{tab:pics}
\scriptsize
\begin{tabular}{l|rrrr|rrrr}
\toprule
 & \multicolumn{4}{c|}{ARM64} & \multicolumn{4}{c}{AMD64} \\
\cmidrule(lr){2-5}\cmidrule(lr){6-9}
Image & MIC & 2-strip & 4-strip & 8-strip & MIC & 2-strip & 4-strip & 8-strip \\
\midrule
MR                   & 649 & 698 & 939 & 1{,}033 & 614 & 281 & 369 & 291\,$^\dagger$ \\
CT                   & 495 & 511 & 987 & 1{,}549 & 444 & 379 & 641 & 736 \\
CR                   & 580 & 865 & 1{,}771 & 3{,}227 & 570 & 768 & 1{,}404 & 2{,}236 \\
XR                   & 668 & 888 & 1{,}728 & 3{,}197 & 626 & 787 & 1{,}408 & 2{,}252 \\
MG1                  & 762 & 1{,}415 & 2{,}403 & 4{,}342 & 916 & 1{,}330 & 1{,}924 & 3{,}575 \\
MG2                  & 753 & 1{,}369 & 2{,}226 & 4{,}109 & 942 & 1{,}302 & 1{,}698 & 3{,}011 \\
MG3                  & 600 & 934 & 1{,}763 & 3{,}158 & 547 & 868 & 1{,}597 & 2{,}781 \\
MG4                  & 788 & 1{,}263 & 2{,}185 & 4{,}120 & 686 & 1{,}258 & 1{,}957 & 3{,}107 \\
CT1                  & 600 & 779 & 1{,}178 & 1{,}639 & 535 & 477 & 743 & 797 \\
CT2                  & 570 & 773 & 1{,}249 & 1{,}296 & 522 & 475 & 718 & 806 \\
MG-N                 & 600 & 988 & 1{,}852 & 3{,}569 & 549 & 871 & 1{,}644 & 2{,}913 \\
MR1                  & 734 & 931 & 1{,}375 & 1{,}898 & 586 & 517 & 729 & 827 \\
MR2                  & 700 & 1{,}003 & 1{,}699 & 2{,}750 & 682 & 786 & 1{,}064 & 1{,}476 \\
MR3                  & 839 & 919 & 1{,}381 & 2{,}035 & 742 & 627 & 861 & 853 \\
MR4                  & 692 & 954 & 1{,}301 & 2{,}063 & 664 & 573 & 813 & 877 \\
NM1                  & 778 & 964 & 1{,}410 & 1{,}868 & 597 & 561 & 777 & 771 \\
RG1                  & 469 & 586 & 1{,}068 & 1{,}964 & 547 & 626 & 1{,}113 & 1{,}833 \\
RG2                  & 652 & 1{,}095 & 1{,}887 & 3{,}269 & 716 & 922 & 1{,}523 & 2{,}494 \\
RG3                  & 685 & 1{,}081 & 2{,}027 & 3{,}289 & 719 & 991 & 1{,}596 & 2{,}504 \\
SC1                  & 699 & 1{,}171 & 2{,}005 & 3{,}174 & 673 & 1{,}042 & 1{,}690 & 2{,}598 \\
XA1                  & 662 & 973 & 1{,}692 & 2{,}404 & 670 & 806 & 1{,}170 & 1{,}443 \\
CT\_BRAIN            & 528 & 627 & 1{,}025 & 1{,}291 & 420 & 426 & 630 & 681 \\
CT\_ANKLE            & 853 & 1{,}082 & 1{,}745 & 1{,}713 & 774 & 654 & 869 & 877 \\
CT\_ORT              & 538 & 679 & 1{,}099 & 1{,}316 & 565 & 485 & 726 & 778 \\
CT\_CHEST            & 500 & 622 & 1{,}042 & 1{,}354 & 446 & 409 & 633 & 691 \\
MR\_HEAD             & 595 & 620 & 915 & 880 & 679 & 306 & 417 & 307 \\
MR\_INTERA           & 590 & 910 & 1{,}575 & 2{,}405 & 454 & 618 & 940 & 1{,}329 \\
DX\_HAND             & 500 & 605 & 1{,}113 & 2{,}051 & 570 & 619 & 1{,}133 & 1{,}697 \\
DX\_CHEST            & 476 & 566 & 1{,}003 & 1{,}971 & 547 & 592 & 990 & 1{,}712 \\
CR\_THORAX           & 552 & 588 & 1{,}128 & 1{,}733 & 561 & 624 & 1{,}141 & 1{,}937 \\
PET1                 & 732 & 637 & 854 & 854 & 679 & 289 & 406 & 317 \\
PET2                 & 656 & 669 & 776 & 923 & 487 & 341 & 356 & 320 \\
CT\_LUNG             & 522 & 689 & 1{,}107 & 1{,}462 & 534 & 468 & 752 & 808 \\
CT\_PANCREAS         & 517 & 506 & 1{,}018 & 1{,}316 & 439 & 390 & 652 & 688 \\
MR\_BRAIN            & 866 & 815 & 999 & 1{,}143 & 708 & 435 & 490 & 448 \\
MR\_BREAST           & 640 & 653 & 920 & 956 & 759 & 288 & 396 & 337 \\
MR\_PROSTATE         & 679 & 742 & 1{,}093 & 1{,}279 & 545 & 321 & 507 & 455 \\
PET\_PSMA$^\ddagger$ & 759 & 727 & -- & -- & 747 & 282 & -- & -- \\
PET\_LUNG$^\ddagger$ & 502 & 359 & -- & -- & 365 & 208 & -- & -- \\
\bottomrule
\end{tabular}

\vspace{2pt}
{\scriptsize $^\dagger$The 130\,KB MR image is too small to benefit from 8-way
parallelism on AMD64: synchronization overhead exceeds the per-strip work, so
4-strip is faster than 8-strip. $^\ddagger$On the two 0.08\,MB PET images,
splitting into four or more strips drives the strip-compressed size to or above
the single-stream size, so their 4- and 8-strip figures reflect near-raw
copying rather than entropy decoding and are omitted.}
\end{table*}

Practical guidance: to decode many small images, use single-threaded MIC
per request, which is already fast. To decode one large image on a multi-core
machine, use strip-parallel mode with four to eight threads.

\subsection{In the browser}

The JavaScript decoder performs well. Notably, the \textbf{four-state}
variant runs 6--21\% \emph{faster} than eight-state in JavaScript, even on the
same hardware where eight states win in C. The reason is the JavaScript engine
(V8): it can keep four independent chains busy but not eight, so the extra
chains add overhead without shortening the critical path. The eight-state
path still works in JavaScript (it lets one decoder accept C-encoded files
without conversion), but four states is the best configuration for the
browser.

Concretely, a 9.35\,MB mammography image decodes in 17.1\,ms (547\,MB/s) using
four-state strips across worker threads. That means a web DICOM viewer can pull
a \texttt{.mic} file straight from object storage and decode it client-side,
with no server doing compression or format conversion. (These numbers are from
Node.js; full in-browser benchmarking is future work.)
Tables~\ref{tab:js-perf} and~\ref{tab:pics-js} give the JS numbers.

\begin{table}[htbp]
\centering
\caption{JavaScript decoder throughput on the ARM64 reference (Apple
M4~Pro, Node.js v22.16), for the four-state and eight-state FSE
variants. Median over 20 iterations after three warm-up runs.}
\label{tab:js-perf}
\begin{tabular}{lrrrrrr}
\toprule
 & & & \multicolumn{2}{c}{4-state} & \multicolumn{2}{c}{8-state} \\
\cmidrule(lr){4-5}\cmidrule(lr){6-7}
Image & Dimensions & Ratio & ms & MB/s & ms & MB/s \\
\midrule
MR\,$^\dagger$ & $256{\times}256$   & 2.35$\times$ &   2.9 &  43 &   2.7 &  46 \\
CT  & $512{\times}512$   & 2.24$\times$ &  12.3 &  41 &  13.1 &  38 \\
CR  & $2140{\times}1760$ & 3.69$\times$ & 137.2 &  52 & 165.5 &  43 \\
MG1 & $2457{\times}1996$ & 8.79$\times$ &  70.4 & 133 &  80.0 & 117 \\
MG3 & $4774{\times}3064$ & 2.29$\times$ & 560.8 &  50 & 607.1 &  46 \\
DX\_HAND & $1480{\times}1410$ & 2.24$\times$ &  82.6 &  48 &  96.1 &  41 \\
PET1 & $256{\times}256$   & 2.74$\times$ &   2.1 &  60 &   2.5 &  50 \\
\bottomrule
\end{tabular}

\vspace{2pt}
{\scriptsize $^\dagger$MR's 0.13\,MB image decodes in under 3\,ms, at the
JavaScript timing floor; the four-state advantage seen on the larger images
is within run-to-run noise here.}
\end{table}

\begin{table}[htbp]
\centering
\caption{PICS parallel JavaScript decoder (\texttt{worker\_threads}) on
the ARM64 reference, with four-state and eight-state FSE per strip.
Worker count was swept over $\{1,2,4,8,14\}$; the best wall-time for
each variant is reported.}
\label{tab:pics-js}
\begin{tabular}{lrrrrr}
\toprule
 & & \multicolumn{2}{c}{4-state strips} & \multicolumn{2}{c}{8-state strips} \\
\cmidrule(lr){3-4}\cmidrule(lr){5-6}
Image & Strips &   ms & MB/s &   ms & MB/s \\
\midrule
MR  & 4 &   0.9 & 134 &   0.9 & 135 \\
CT  & 4 &   4.1 & 123 &   4.5 & 112 \\
CR  & 8 &  22.9 & 313 &  28.0 & 257 \\
MG1 & 8 &  17.1 & 547 &  16.7 & 559 \\
DX\_HAND & 8 &  17.1 & 233 &  17.8 & 224 \\
\bottomrule
\end{tabular}
\end{table}

% ============================================================
\section{Discussion}
\label{sec:discussion}

\subsection{Reading the results together}

Three observations tie the numbers together. First, MIC's \textbf{ratio
advantage over Zstandard} (about 10\% in geometric mean, higher on 34 of 39
images) comes mainly from working on 16-bit residuals directly instead of
splitting them into bytes, and that same choice removes the per-byte branches
that slow decoding. Second, on the largest mammography images on AMD64,
\textbf{HTJ2K's block coder} regains the speed lead by amortizing its setup
over many pixels; MIC stays within 0--20\% there.
Third, \textbf{multi-state decoding} is what makes MIC's decoder competitive
with HTJ2K's vectorized block coder, and it costs nothing in file size.

\subsection{Choosing a codec}

For \textbf{low-latency viewing} (PACS, thumbnails, tomosynthesis playback),
single-threaded MIC is a sensible default (${\sim}600$--900\,MB/s per request),
with strip-parallel MIC for one large image at a time. For \textbf{smallest
archival files}, JPEG-LS keeps a ${\sim}14\%$ better ratio, at the cost of
decoding 1.6--5.7$\times$ slower on ARM64. Where \textbf{DICOM transfer-syntax
compatibility} is required, HTJ2K is the appropriate choice, since it is a
standard and is particularly strong on large mammography images on AMD64. Where
\textbf{no native code is permitted} (browsers, sandboxes), the pure-Go
reference or the 20\,KB JavaScript decoder runs at about half the C decoder's
speed.

A more structured comparison---a Pareto view of ratio vs.\ decode
speed~\cite{geldreich_pareto,jpegxl_pareto}, plus a weighted decision matrix
over ratio, speed, platform coverage, and deployment
ease~\cite{saaty_ahp,sayood}, following deployment taxonomies for DICOM
compression~\cite{clunie_miit2009,liu2017}---puts MIC, HTJ2K, and JPEG-LS on
the efficiency frontier (MIC for throughput, JPEG-LS for ratio, HTJ2K in
between), with Delta + Zstandard dominated by MIC on both axes. HTJ2K remains
the choice whenever standard compatibility is a hard requirement that the
scoring does not capture.
Table~\ref{tab:decision-matrix} reports the weighted matrix.

\begin{table}[htbp]
\centering
\caption{Codec selection matrix. Each criterion is normalised so the
best codec on that row scores~1.00. Quantitative rows are geometric
means across the 39-image set. Browser compatibility: shipped
purpose-built JS or WASM decoder used in production~=~1.00; WASM port of
the native library used in clinical viewers~=~0.85; no browser
path~=~0.0. Deployment ease: zero external runtime dependencies~=~1.00;
ubiquitous system dependency~=~0.90; C/C++ library requiring CMake build
and link integration~=~0.70. Weight columns $w_P$, $w_A$, $w_B$ reflect
typical priorities for (P)~PACS viewer / low-latency decode, (A)~archival
storage, (B)~browser delivery.}
\label{tab:decision-matrix}
\scriptsize
\setlength{\tabcolsep}{3pt}
\begin{tabular}{lrrrrrr}
\toprule
Criterion & $w_P$ & $w_A$ & $w_B$ & MIC & HTJ2K & JPEG-LS \\
\midrule
Compression ratio                & 0.20 & 0.50 & 0.10 & 0.88 & 0.89 & \textbf{1.00} \\
Decompression throughput (ARM64) & 0.30 & 0.10 & 0.15 & \textbf{1.00} & 0.57 & 0.34 \\
Encoding throughput (ARM64)      & 0.05 & 0.10 & 0.05 & \textbf{1.00} & 0.40 & 0.25 \\
AMD64 native                     & 0.10 & 0.05 & 0.05 & \textbf{1.00} & \textbf{1.00} & \textbf{1.00} \\
ARM64 native                     & 0.10 & 0.05 & 0.05 & \textbf{1.00} & \textbf{1.00} & \textbf{1.00} \\
Browser compatibility            & 0.05 & 0.00 & 0.40 & \textbf{1.00} & 0.85 & 0.85 \\
Deployment ease                  & 0.20 & 0.20 & 0.20 & \textbf{1.00} & 0.70 & 0.70 \\
\midrule
\textbf{Weighted score (P) viewer}   & --- & --- & --- & \textbf{0.98} & 0.75 & 0.70 \\
\textbf{Weighted score (A) archival} & --- & --- & --- & \textbf{0.94} & 0.78 & 0.80 \\
\textbf{Weighted score (B) browser}  & --- & --- & --- & \textbf{0.99} & 0.77 & 0.74 \\
\bottomrule
\end{tabular}
\end{table}

\subsection{Limitations}

\textbf{Small dataset.} 39 images across eleven modalities is varied but far from
a clinical archive; broader validation is needed. \textbf{Tuning on the test
set.} The adaptive table-size thresholds and the predictor choice were informed
by this dataset; the rules are simple (two thresholds, one predictor), which
limits overfitting risk, but larger studies should confirm they generalize.
\textbf{Baseline coverage.} Stronger byte-oriented baselines (Delta +
byte/bit-shuffle + Zstd, DICOM RLE) would more sharply isolate the value of the
16-bit entropy stage; left to future work. \textbf{Reporting.} Tables give
single representative throughput values rather than full confidence intervals.
\textbf{Not yet a DICOM transfer syntax.} MIC works on the extracted pixel
array; real PACS integration would need a private encapsulation or a
standardization step.

% ============================================================
\section{Conclusion}
\label{sec:conclusion}

MIC is a 16-bit-native lossless codec for medical images: spatial prediction,
16-bit run-length coding, a large-alphabet table-based entropy coder, and a
multi-state decoder for speed. On 39 DICOM images it reaches a geometric-mean
ratio of 3.36$\times$ (within ${\sim}2\%$ of HTJ2K, ${\sim}88\%$ of JPEG-LS)
while being the fastest decoder on all 39 images on ARM64 and 36/39 on AMD64,
and the fastest encoder on 38/39 ARM64 images and 32/39 on AMD64. A 20\,KB JavaScript
decoder and a WebAssembly build show the format can be decoded client-side in a
browser. Future work: larger multi-hospital datasets; stronger shuffle-based
baselines; full statistical reporting; an ARM64 vector kernel for the wavelet
variant; and benchmarking the JavaScript decoder in real browsers. MIC is open
source: \url{https://github.com/pappuks/medical-image-codec}.

% ============================================================
\begin{thebibliography}{99}

\bibitem{dicom}
NEMA, ``Digital Imaging and Communications in Medicine (DICOM),''
PS3.5---Data Structures and Encoding,
\url{https://www.dicomstandard.org/}, 2024.

\bibitem{taubman2002}
D.\,S.\,Taubman and M.\,W.\,Marcellin,
\textit{JPEG2000: Image Compression Fundamentals, Standards and Practice}.
Springer, 2002.

\bibitem{htj2k2019}
D.\,S.\,Taubman, ``High throughput JPEG 2000 (HTJ2K): Algorithm, performance
evaluation, and potential,'' \textit{Proc.\ SPIE}, vol.\,11137, 2019.

\bibitem{taubman2022}
D.\,Taubman, A.\,Naman, M.\,Smith, P.-A.\,Lemieux, H.\,Saadat, O.\,Watanabe,
and R.\,Mathew,
``High Throughput JPEG~2000 for Video Content Production and Delivery Over IP
Networks,''
\textit{Front.\ Signal Process.}, vol.\,2, Article 885644, Apr.\,2022.

\bibitem{loco1}
M.\,J.\,Weinberger, G.\,Seroussi, and G.\,Sapiro,
``The LOCO-I lossless image compression algorithm: Principles and
standardization into JPEG-LS,''
\textit{IEEE Trans.\ Image Process.}, vol.\,9, no.\,8, pp.\,1309--1324, 2000.

\bibitem{wu1997}
X.\,Wu and N.\,Memon,
``Context-based, adaptive, lossless image coding,''
\textit{IEEE Trans.\ Commun.}, vol.\,45, no.\,4, pp.\,437--444, 1997.

\bibitem{sneyers2016}
J.\,Sneyers and P.\,Wuille,
``FLIF: Free Lossless Image Format based on MANIAC Compression,''
\textit{Proc.\ IEEE ICIP}, 2016.

\bibitem{duda2009}
J.\,Duda, ``Asymmetric numeral systems,'' arXiv:0902.0271, 2009.

\bibitem{duda2013}
J.\,Duda, ``Asymmetric numeral systems: entropy coding combining speed of
Huffman coding with compression rate of arithmetic coding,''
arXiv:1311.2540, 2013.

\bibitem{fse}
Y.\,Collet, ``Finite State Entropy---A new breed of entropy coder,''
\url{https://github.com/Cyan4973/FiniteStateEntropy}, 2013.

\bibitem{rfc8878}
Y.\,Collet and M.\,Kucherawy,
``Zstandard Compression and the `application/zstd' Media Type,''
RFC~8878, IETF, Feb.\,2021.

\bibitem{giesen2014}
F.\,Giesen, ``Interleaved entropy coders,'' arXiv:1402.3392, 2014.

\bibitem{lin2023}
F.\,Lin, K.\,Arunruangsirilert, H.\,Sun, and J.\,Katto,
``Recoil: Parallel rANS Decoding with Decoder-Adaptive Scalability,''
\textit{Proc.\ 52nd ICPP '23}, pp.\,31--40, 2023.

\bibitem{weissenberger2019}
A.\,Weissenberger and B.\,Schmidt,
``Massively Parallel ANS Decoding on GPUs,''
\textit{Proc.\ 48th ICPP '19}, Article~100, 2019.

\bibitem{legall1988}
D.\,Le Gall and A.\,Tabatabai,
``Sub-band coding of digital images using symmetric short kernel filters and
arithmetic coding techniques,'' \textit{Proc.\ IEEE ICASSP}, pp.\,761--764,
1988.

\bibitem{png}
W3C, ``Portable Network Graphics (PNG) Specification (Second Edition),''
ISO/IEC 15948:2003, Nov.\,2003.

\bibitem{williams2009}
S.\,Williams, A.\,Waterman, and D.\,Patterson,
``Roofline: An insightful visual performance model for multicore
architectures,''
\textit{Commun.\ ACM}, vol.\,52, no.\,4, pp.\,65--76, Apr.\,2009.

\bibitem{liu2017}
F.\,Liu \textit{et al.},
``The Current Role of Image Compression Standards in Medical Imaging,''
\textit{Information}, vol.\,8, no.\,4, p.\,131, Oct.\,2017.

\bibitem{clunie2000}
D.\,A.\,Clunie,
``Lossless Compression of Grayscale Medical Images: Effectiveness of
Traditional and State-of-the-Art Approaches,''
\textit{Proc.\ SPIE Medical Imaging}, 2000.

\bibitem{openjph}
A.\,N.\,Aous,
``OpenJPH: An open-source implementation of High-Throughput JPEG~2000,''
\url{https://github.com/aous72/OpenJPH}, 2024.

\bibitem{charls}
Team CharLS,
``CharLS: A C/C++ JPEG-LS library implementation,''
\url{https://github.com/team-charls/charls}, 2024.

\bibitem{nema_wg04}
NEMA Working Group~4 (WG-04),
``DICOM Compression Sample Images,''
National Electrical Manufacturers Association,
\url{https://www.dicomstandard.org/Resources/Open-Content/Compression-Samples},
2024.

\bibitem{clunie_tomo}
D.\,A.\,Clunie,
``Breast Tomosynthesis Case~1,'' DICOM sample images,
\url{http://www.dclunie.com/pixelmedimagecases/tomo/case1/}, 2009.

\bibitem{geldreich_pareto}
R.\,Geldreich, ``Quick survey of the lossless decompression Pareto frontier,''
2015, \url{http://richg42.blogspot.com/2015/11/the-lossless-decompression-pareto.html}.

\bibitem{jpegxl_pareto}
J.\,Sneyers, ``JPEG XL and the Pareto front,'' Cloudinary Blog, 2022,
\url{https://cloudinary.com/blog/jpeg-xl-and-the-pareto-front}.

\bibitem{saaty_ahp}
T.\,L.\,Saaty, \textit{The Analytic Hierarchy Process: Planning, Priority
Setting, Resource Allocation}. New York, NY, USA: McGraw-Hill, 1980.

\bibitem{sayood}
K.\,Sayood, \textit{Introduction to Data Compression}, 5th ed.
Burlington, MA, USA: Morgan Kaufmann, 2017.

\bibitem{clunie_miit2009}
D.\,A.\,Clunie, ``DICOM support for compression: More than JPEG,''
\textit{Medical Imaging Informatics and Teleradiology (MIIT)}, 2009,
\url{https://www.dclunie.com/papers/MIIT_2009_DicomCompression_Clunie.pdf}.

\end{thebibliography}

\end{document}
